\chapter{Von zeitlicher Exposition zur modellierten Antwort}
\label{chap:runtime-dynamics}

Das Durchfahrtsgesetz exponiert das bewegte System nun gegenüber
\(\kappa(s(t))\), \(u(s(t))\), dem Profil und dem qualifizierten Gleis-Frame.
Dieses Kapitel markiert die Grenze zwischen Größen, die aus dieser Kinematik
folgen, und Größen, die ein Fahrzeug--Gleis-Antwortmodell benötigen.

\begin{principle}[Laufzeitdynamikprinzip]
\label{princ:runtime-dynamics}
Laufzeitdynamik konsumiert einen qualifizierten Eisenbahnzustand und ein
Durchfahrtsgesetz. Jede behauptete Antwort muss zusätzlich das erzeugende
Modell identifizieren.
\end{principle}

\section{Ein Ergebnis, keine zweite Zustandsdefinition}

\begin{definition}[Dynamischer Laufzeitzustand]
\label{def:runtime-dynamic-state}
Für den engen Zweck dieses Teils ist ein dynamischer Laufzeitzustand die
Ausgabe eines deklarierten Antwortmodells, ausgewertet für qualifizierten
Eisenbahnzustand, Durchfahrt und Modelleingaben.
\end{definition}

Diese Definition klassifiziert eine Ausgabe; sie definiert weder pose3,
physischen Zustand noch Beobachtung neu. Eine vorhergesagte
Wagenkastenbeschleunigung, eine beobachtete Radlast und ein reduzierter
kinematischer Indikator sind verschiedene Ergebnisarten, auch wenn sie sich
auf dasselbe \(s(t)\) beziehen.

\begin{definition}[Laufzeitdynamikoperator]
\label{def:runtime-dynamics-operator}
\[
\mathcal D_M:
\bigl(\mathcal S_{\mathrm{rail}}(s(t)),s(t),Q_M\bigr)
\longmapsto x_M(t)
\]
bezeichne die Auswertung durch ein explizit identifiziertes Antwortmodell
\(M\) mit qualifizierten Modelleingaben \(Q_M\).
\end{definition}

Der Modellindex ist wesentlich. Es gibt kein universelles \(\mathcal D\), das
pose3 stillschweigend in validiertes Fahrzeugverhalten umwandelt.

\begin{definition}[Laufzeitzustandsentwicklung]
\label{def:runtime-dynamics-state-evolution}
Laufzeitzustandsentwicklung ist die zeit- oder raumindizierte
Ausgabetrajektorie \(x_M\), die das gewählte Modell über dem anwendbaren
Durchfahrtsintervall erzeugt.
\end{definition}

\section{Reduzierte Kinematik}

Unter dem rechtshändigen Gleis-Frame aus Teil V, seinem angegebenen
Überhöhungsvorzeichen, starrer punktförmiger Folgebewegung und unter
Vernachlässigung von Federung, Kontakt, Gleisnachgiebigkeit, Unregelmäßigkeiten,
Wind und weiteren Effekten ist
\[
a_\kappa(t)=v(t)^2\kappa(s(t))
\]
eine kinematische Querkomponente. Unter der angegebenen
Überhöhungswinkelkonvention ist
\[
a_{\mathrm{eff}}(t)
=v(t)^2\kappa(s(t))-g\sin\phi(s(t))
\]
ein reduzierter erklärender Indikator. Bei veränderlicher Geschwindigkeit gilt
\begin{align}
\dot a_{\mathrm{eff}}
={}&2v\dot v\,\kappa
  +v^3\kappa'
  -gv\cos\phi\,\phi',
\label{eq:variable-speed-effective-rate}
\end{align}
wobei alle räumlichen Größen an \(s(t)\) ausgewertet werden. Der Term
\(2v\dot v\kappa\) fehlt in der Konstantgeschwindigkeitsbeziehung aus Teil V.
Dieser Indikator ist für sich weder Komfort, Sicherheit, Verschleiß,
Radentlastung noch Fahrzeugantwort.

\section{Die explizite Modellgrenze}

\begin{center}
\fbox{\begin{minipage}{0.88\textwidth}
\textbf{Hier endet die Kinematik.} Eisenbahnzustand und Durchfahrt können
Trajektorie, Frame, Krümmungs- und Überhöhungsexposition sowie reduzierte
Beschleunigungs- oder Ratenindikatoren unter angegebenen Annahmen liefern.

\medskip
\textbf{Vorhergesagte Antwort beginnt erst, wenn ein Modell ergänzt:}
Fahrzeugmasse und Trägheit, Federung, Rad--Schiene-Kontakt, Gleissteifigkeit
und -dämpfung, Unregelmäßigkeiten, Betriebslast, Umweltzustand und relevante
Instandhaltungsbedingung. Modellanwendbarkeit, Kalibrierung und Unsicherheit
bleiben an jeder Ausgabe erhalten.
\end{minipage}}
\end{center}

Nicht jede Verwendung benötigt jede aufgezählte Eingabe. Eine Auslassung muss
jedoch aus dem angegebenen Geltungsbereich des Modells folgen und darf nicht
stillschweigend verschwinden. Eine gemessene Antwort bleibt Beobachtung und
wird nicht durch Benennung zur Modellausgabe oder zum physischen Zustand.

\begin{principle}[Mechaniktrennungsprinzip]
\label{princ:mechanics-separation}
Reduzierte Eisenbahnkinematik und eine vollständige
Fahrzeug--Gleis-Mechaniksimulation sind verschiedene Operationen mit
unterschiedlichen Eingaben, Anwendbarkeiten und Evidenzen.
\end{principle}

\section{Der laufende Vergleich}

Die beiden Fälle aus Teil V verwenden identische \(\kappa\) und Überhöhung,
aber unterschiedliche konstante Testgeschwindigkeiten; daher unterscheiden
sich ihre reduzierten \(v^2\kappa\)-Terme. Das Beispiel mit veränderlicher
Geschwindigkeit ergänzt eine von \(\dot v\) abhängige Exposition. Ein anderes
\(u(s)\) würde Gleis-Frame und reduzierte effektive Komponente ändern, nicht
die Identität von \(\kappa(s)\). Selbst identische reduzierte Indikatoren
garantieren dennoch keine identische Antwort: Fahrzeuge, Gleiszustand, Last und
Umwelt können verschieden sein.

\section{AXTRAN und deklarierte Ziele}

Der mit axtranNew/AXTRAN verbundene Alignment-Rechenkern kann Geometrie oder
freie Parameter vorschlagen. Eine deklarierte reduzierte Größe kann später als
Ziel oder Nebenbedingung verwendet werden, wenn Annahmen, Zweck und
Anwendbarkeit explizit sind. Das Ziel ist eine der Berechnung zugeführte
Engineering-Wahl und wird vom Solver nicht als Autorität entdeckt. Ein
berechnetes Optimum ist weder automatisch anwendbar noch angenommen.

Die gegenwärtige Implementierung begründet keine allgemeine gemeinsame
Optimierung von Geometrie, Überhöhung, Durchfahrt und Fahrzeugantwort. Der neue
Kontinuitätssolver wird daher nicht als Dynamikoptimierer beschrieben.

\begin{summary}
Zeitliche Exposition folgt aus \(s(t)\). Reduzierte Kinematik folgt aus
angegebenen Frame- und Bewegungsannahmen. Fahrzeug--Gleis-Antwort beginnt erst
mit einem identifizierten Modell. Bewertung und Entscheidung bleiben
nachgelagert.
\end{summary}
